\begin{tikzpicture}[scale=0.85, every node/.style={font=\small}]
	
	% ============================================================
	% 1. ARDUINO NANO 33 BLE SENSE (aus arduinonano33bletikz.tex)
	% ============================================================
	\begin{scope}[xshift=0cm, yshift=0cm]
		% Hauptplatine - Farbe c0f7391 (15,115,145)
		\definecolor{arduinoBlue}{RGB}{15,115,145}
		\draw[fill=arduinoBlue, draw=arduinoBlue!30!black, thick, rounded corners=2pt] 
		(0,0) rectangle (4.5,6.5);
		
		% Name
		\node[text=white, font=\bfseries\small] at (2.25,6.1) {Arduino Nano 33 BLE Sense};
		
		% USB-Buchse (vereinfacht)
		\draw[fill=gray!40, draw=gray!60, thick, rounded corners=1pt] 
		(1.8,6.2) rectangle (2.7,6.5);
		
		% Linke Pins (D13-D0) - Farbe c333333 (51,51,51)
		\definecolor{pinColor}{RGB}{51,51,51}
		\foreach \i/\label in {0/D13, 1/D12, 2/D11, 3/D10, 4/D9, 5/D8, 
			6/D7, 7/D6, 8/D5, 9/D4, 10/D3, 11/D2, 12/GND, 13/RESET, 14/RX}
		{
			\pgfmathsetmacro{\y}{5.7 - \i*0.35}
			\draw[fill=cf7bd13, draw=black!50] (0.1,\y-0.1) rectangle (0.3,\y+0.1);
			\node[text=white, font=\tiny, anchor=east] at (0.05,\y) {\label};
		}
		
		% Rechte Pins (3V3-TX) - Farbe c333333
		\foreach \i/\label in {0/+3V3, 1/AREF, 2/A0, 3/A1, 4/A2, 5/A3, 
			6/A4, 7/A5, 8/A6, 9/A7, 10/+5V, 11/GND, 12/VIN, 13/RESET, 14/TX}
		{
			\pgfmathsetmacro{\y}{5.7 - \i*0.35}
			\draw[fill=cf7bd13, draw=black!50] (4.2,\y-0.1) rectangle (4.4,\y+0.1);
			\node[text=white, font=\tiny, anchor=west] at (4.45,\y) {\label};
		}
		
		% Prozessor (nRF52840)
		\draw[fill=black!80, draw=black, thick, rounded corners=1pt] 
		(1.5,2.5) rectangle (3.0,4.5);
		\node[text=white, font=\tiny] at (2.25,3.5) {nRF52840};
		
		% LEDs
		\draw[fill=green!60!black, draw=black] (1.2,1.2) circle (0.12);
		\draw[fill=red!60!black, draw=black] (1.6,1.2) circle (0.12);
		\draw[fill=orange!60!black, draw=black] (2.0,1.2) circle (0.12);
		
		% Arduino Logo (vereinfacht)
		\node[text=white, font=\tiny] at (2.25,0.5) {ARDUINO\textregistered};
	\end{scope}
	
	% ============================================================
	% 2. SERVO (SG90) - aus AktorBild.tex
	% ============================================================
	\begin{scope}[xshift=7cm, yshift=3.5cm]
		% Hauptgehäuse - Farbe blue!70!black
		\draw[fill=blue!70!black, draw=blue!30!black, thick, rounded corners=2pt]
		(0,0) rectangle (2.0,2.2);
		
		% Transparente Innenfläche
		\draw[fill=blue!45, draw=blue!20!black, opacity=0.35]
		(0.15,0.15) rectangle (1.85,2.05);
		
		% Befestigungslaschen oben
		\draw[fill=blue!75!black, draw=blue!30!black, thick, rounded corners=1pt]
		(-0.3,1.6) rectangle (0,1.8);
		\draw[fill=blue!75!black, draw=blue!30!black, thick, rounded corners=1pt]
		(2.0,1.6) rectangle (2.3,1.8);
		
		% Oberes Getriebegehäuse
		\draw[fill=blue!80!black, draw=blue!30!black, thick, rounded corners=2pt]
		(0.6,2.2) rectangle (1.4,2.4);
		
		% Servoarm
		\draw[fill=gray!60, draw=gray!40!black, thick, rounded corners=2pt]
		(0.2,2.4) rectangle (1.8,2.55);
		
		% Zentrale Verdickung
		\draw[fill=gray!70, draw=gray!40!black, thick, rounded corners=1pt]
		(0.9,2.38) rectangle (1.3,2.57);
		
		% Frontlabel - Farbe black!85, yellow!70!black
		\draw[fill=black!85, draw=yellow!70!black, line width=0.7pt]
		(0.4,0.5) rectangle (1.6,1.4);
		\draw[draw=yellow!80!black, line width=0.4pt]
		(0.45,0.55) rectangle (1.55,1.35);
		\node[text=yellow!80!white, font=\tiny\bfseries] at (1.0,1.1) {TowerPro};
		\node[text=white, font=\tiny] at (1.0,0.85) {Micro Servo};
		\node[text=white, font=\tiny] at (1.0,0.65) {SG90};
		
		% Gehäusekanten
		\draw[blue!20!white, line width=0.5pt, opacity=0.75]
		(0.15,0.15) -- (0.15,2.05);
		\draw[blue!20!white, line width=0.5pt, opacity=0.75]
		(1.85,0.15) -- (1.85,2.05);
		
		% Kabelausgang links unten
		\draw[fill=black!85, draw=black, thick, rounded corners=1pt]
		(0,0.2) rectangle (-0.25,0.7);
		
		% Kabel (Orange=Signal, Rot=VCC, Braun=GND)
		\draw[brown, line width=2pt] (-0.25,0.3) -- (-1.5,0.3);
		\draw[red, line width=2pt] (-0.25,0.45) -- (-1.5,0.45);
		\draw[orange, line width=2pt] (-0.25,0.6) -- (-1.5,0.6);
		
		% Stecker links
		\draw[fill=black!85, draw=black, thick, rounded corners=1pt]
		(-2.0,0.1) rectangle (-1.5,0.8);
		
		% Kontakte am Stecker
		\draw[fill=yellow!70!white, draw=yellow!50!black] (-1.85,0.25) circle (0.05);
		\draw[fill=yellow!70!white, draw=yellow!50!black] (-1.85,0.45) circle (0.05);
		\draw[fill=yellow!70!white, draw=yellow!50!black] (-1.85,0.65) circle (0.05);
		
		% Kabelbeschriftung
		\node[font=\tiny, anchor=east] at (-2.1,0.25) {GND};
		\node[font=\tiny, anchor=east] at (-2.1,0.45) {VCC};
		\node[font=\tiny, anchor=east] at (-2.1,0.65) {Signal};
		endscope
		
		% ============================================================
		% 3. ULTRASCHALLSENSOR (HC-SR04) - aus TikzUltraschallsensor.tex
		% ============================================================
		\begin{scope}[xshift=6.5cm, yshift=-1.5cm, scale=0.9]
			% Farben aus der Vorlage
			\definecolor{pcbblue}{RGB}{0,50,120}
			\definecolor{sensorbody}{RGB}{180,180,180}
			\definecolor{sensorinner}{RGB}{220,220,220}
			\definecolor{pincolor}{RGB}{247,189,19}
			
			% Hauptplatine
			\draw[fill=pcbblue, rounded corners=2pt, draw=black!70] 
			(0,0) rectangle (5.0,2.5);
			
			% Befestigungslöcher
			\foreach \x/\y in {0.25/0.25, 4.75/0.25, 0.25/2.25, 4.75/2.25} {
				\draw[fill=white, draw=gray] (\x,\y) circle (0.12);
			}
			
			% Quarz
			\draw[fill=gray!40, rounded corners=3pt] (1.8,2.1) rectangle (3.2,2.35);
			
			% Produktname
			\node[text=white, font=\sffamily\bfseries\small] at (2.5,1.6) {HC - SR04};
			
			% Sender & Empfänger
			\foreach \xpos in {1.1, 3.9} {
				% Äußerer Ring
				\draw[fill=white, draw=gray] (\xpos,1.0) circle (0.85);
				% Metallgehäuse
				\draw[fill=sensorbody] (\xpos,1.0) circle (0.75);
				% Gitterbereich
				\draw[fill=sensorinner] (\xpos,1.0) circle (0.6);
				% Gitter-Muster
				\begin{scope}
					\clip (\xpos,1.0) circle (0.6);
					\draw[step=0.06, black!30, thin] (\xpos-0.6,1.0-0.6) grid (\xpos+0.6,1.0+0.6);
				\end{scope}
				% Innerster Kreis
				\draw[fill=black!60] (\xpos,1.0) circle (0.3);
			}
			
			% Beschriftung
			\node[text=white, font=\tiny, anchor=north] at (1.1,0.05) {Sender};
			\node[text=white, font=\tiny, anchor=north] at (3.9,0.05) {Empfänger};
			
			% Pins
			\foreach \i/\txt in {0/Vcc, 1/Trig, 2/Echo, 3/Gnd} {
				\draw[fill=gray!40, draw=black!50] (2.0 + \i*0.35, 0.08) rectangle (2.1 + \i*0.35, -0.25);
				\draw[fill=pincolor] (1.9 + \i*0.35, 0.05) rectangle (2.2 + \i*0.35, 0.3);
				\node[text=white, font=\tiny, rotate=90, anchor=west] at (2.05 + \i*0.35, 0.35) {\txt};
			}
			endscope
			
			% ============================================================
			% 4. VERBINDUNGEN
			% ============================================================
			
			% --- Verbindung: Arduino +5V -> Servo VCC (rot) ---
			\draw[red, line width=2pt] (4.45,1.75) -- (5.2,1.75) -- (5.2,5.45) -- (5.5,5.45);
			\node[text=red, font=\tiny, anchor=south] at (5.2,5.55) {+5V};
			
			% --- Verbindung: Arduino GND (links) -> Servo GND (braun) ---
			\draw[brown, line width=2pt] (0.25,2.35) -- (0.25,6.0) -- (5.5,6.0);
			\node[text=brown, font=\tiny, anchor=south] at (2.5,6.1) {GND};
			
			% --- Verbindung: Arduino D9 (PWM) -> Servo Signal (orange) ---
			\draw[orange, line width=2pt] (0.25,3.75) -- (0.25,5.7) -- (5.5,5.7);
			\node[text=orange, font=\tiny, anchor=south] at (2.0,5.8) {D9 (PWM)};
			
			% --- Verbindung: Arduino +5V -> HC-SR04 Vcc ---
			\draw[red!70!black, line width=2pt, dashed] (4.45,1.2) -- (5.8,1.2) -- (5.8,0.3) -- (6.2,0.3);
			\node[text=red!70!black, font=\tiny, anchor=south] at (5.0,1.3) {+5V};
			
			% --- Verbindung: Arduino GND (rechts) -> HC-SR04 Gnd ---
			\draw[brown!70!black, line width=2pt, dashed] (4.45,2.0) -- (5.8,2.0) -- (5.8,-0.1) -- (6.2,-0.1);
			\node[text=brown!70!black, font=\tiny, anchor=south] at (5.0,2.1) {GND};
			
			% --- Verbindung: Arduino D10 -> HC-SR04 Trig ---
			\draw[blue!60!white, line width=2pt, dashed] (0.25,4.1) -- (-0.5,4.1) -- (-0.5,-0.1) -- (6.1,-0.1);
			\node[text=blue!60!white, font=\tiny, anchor=east] at (-0.6,4.1) {D10 (Trig)};
			
			% --- Verbindung: Arduino D11 -> HC-SR04 Echo ---
			\draw[green!60!black, line width=2pt, dashed] (0.25,4.45) -- (-0.5,4.45) -- (-0.5,0.1) -- (6.1,0.1);
			\node[text=green!60!black, font=\tiny, anchor=east] at (-0.6,4.45) {D11 (Echo)};
			
			% ============================================================
			% 5. LEGENDE mit den originalen Farben
			% ============================================================
			\begin{scope}[xshift=0cm, yshift=-3.0cm]
				\draw[fill=white!90!black, draw=black!50, rounded corners=2pt] 
				(0,0) rectangle (12,1.8);
				
				\node[font=\bfseries] at (0.5,1.4) {Legende:};
				
				% Servo - mit den Farben aus AktorBild.tex
				\draw[red, line width=2pt] (1.5,1.3) -- (2.0,1.3);
				\node[anchor=west] at (2.1,1.3) {VCC (+5V)};
				
				\draw[brown, line width=2pt] (3.8,1.3) -- (4.3,1.3);
				\node[anchor=west] at (4.4,1.3) {GND};
				
				\draw[orange, line width=2pt] (6.1,1.3) -- (6.6,1.3);
				\node[anchor=west] at (6.7,1.3) {Signal (PWM)};
				
				% HC-SR04 - mit den Farben aus TikzUltraschallsensor.tex
				\draw[red!70!black, line width=2pt, dashed] (1.5,0.8) -- (2.0,0.8);
				\node[anchor=west] at (2.1,0.8) {+5V};
				
				\draw[brown!70!black, line width=2pt, dashed] (3.8,0.8) -- (4.3,0.8);
				\node[anchor=west] at (4.4,0.8) {GND};
				
				\draw[blue!60!white, line width=2pt, dashed] (6.1,0.8) -- (6.6,0.8);
				\node[anchor=west] at (6.7,0.8) {Trig (D10)};
				
				\draw[green!60!black, line width=2pt, dashed] (8.5,0.8) -- (9.0,0.8);
				\node[anchor=west] at (9.1,0.8) {Echo (D11)};
				
				% Zusätzliche Info
				\node[font=\tiny, text=gray] at (6.0,0.3) {Verbindungen: durchgezogen = Servo, gestrichelt = HC-SR04};
				endscope
				
			\end{tikzpicture}